% ============================================================
\chapter{Environnements et Gestion de D\'ependances}
\label{ch:environnements}
\index{environnements}
% ============================================================

\section{Pourquoi isoler les environnements~?}
\label{sec:why-envs}

Un probl\`eme classique en ML~: <<~\c{C}a marche sur ma machine~!~>>
(\emph{works on my machine}). Les causes~:
\begin{itemize}
  \item Versions de Python diff\'erentes
  \item Versions de biblioth\`eques incompatibles
  \item D\'ependances syst\`eme manquantes (CUDA, cuDNN)
  \item Variables d'environnement non configur\'ees
\end{itemize}

\begin{definition}[Environnement virtuel]
Un \textbf{environnement virtuel}\index{environnement virtuel} est un
r\'epertoire isol\'e contenant une installation Python sp\'ecifique et un
ensemble de packages, ind\'ependant du syst\`eme et des autres projets.
\end{definition}

\begin{center}
\begin{tikzpicture}[
  env/.style={rectangle, draw, rounded corners, minimum width=3cm,
    minimum height=2.5cm, align=center, font=\small},
  pkg/.style={rectangle, draw, fill=blue!10, font=\scriptsize,
    minimum width=2.5cm, rounded corners}
]
  \node[env, fill=green!5] (env1) {};
  \node[above=0.1cm] at (env1.north) {\textbf{Projet A}};
  \node[pkg] at ([yshift=0.3cm]env1.center) {scikit-learn 1.3};
  \node[pkg] at ([yshift=-0.3cm]env1.center) {pandas 2.0};
  \node[below=0.1cm, font=\scriptsize] at (env1.south) {Python 3.10};

  \node[env, fill=orange!5, right=2cm of env1] (env2) {};
  \node[above=0.1cm] at (env2.north) {\textbf{Projet B}};
  \node[pkg] at ([yshift=0.3cm]env2.center) {scikit-learn 1.5};
  \node[pkg] at ([yshift=-0.3cm]env2.center) {pandas 2.2};
  \node[below=0.1cm, font=\scriptsize] at (env2.south) {Python 3.12};

  \node[below=1.5cm of env1, font=\small] (sys) {Syst\`eme};
  \node[below=1.5cm of env2, font=\small] (sys2) {Syst\`eme};
  \draw[dashed, gray] ([xshift=-1cm]sys.north) -- ([xshift=1cm]sys2.north);
\end{tikzpicture}
\end{center}

\section{venv --- L'outil int\'egr\'e \`a Python}
\label{sec:venv}
\index{venv}

\cli{venv} est le module standard de Python pour cr\'eer des environnements
virtuels. Il est simple, l\'eger, et ne n\'ecessite aucune installation
suppl\'ementaire.

\begin{codeblock}[title=\textbf{Commandes venv essentielles}]
\begin{minted}{bash}
# Creer un environnement virtuel
python -m venv .venv

# Activer (Linux/macOS)
source .venv/bin/activate

# Activer (Windows)
.venv\Scripts\activate

# Installer des packages
pip install numpy pandas scikit-learn

# Sauvegarder les dependances
pip freeze > requirements.txt

# Recreer l'environnement
pip install -r requirements.txt

# Desactiver
deactivate
\end{minted}
\end{codeblock}

\begin{attention}[title=\textbf{Pi\`eges de pip freeze}]
\cli{pip freeze} inclut \textbf{toutes} les d\'ependances transitives, ce
qui rend le fichier fragile. Pr\'ef\'erez~:
\begin{itemize}
  \item Maintenir manuellement un \file{requirements.txt} avec les
    d\'ependances directes et leurs versions.
  \item Utiliser \pkg{pip-tools} pour g\'en\'erer un \file{requirements.txt}
    verrouill\'e \`a partir d'un \file{requirements.in}.
\end{itemize}
\end{attention}

\subsection{pip-tools~: le verrouillage des d\'ependances}
\index{pip-tools}

\begin{codeblock}[title=\textbf{Workflow pip-tools}]
\begin{minted}{bash}
# Installer pip-tools
pip install pip-tools

# requirements.in (dependances directes seulement)
# numpy>=1.24
# pandas>=2.0
# scikit-learn>=1.3

# Compiler (resoudre et verrouiller)
pip-compile requirements.in -o requirements.txt

# Installer exactement les versions verrouillees
pip-sync requirements.txt
\end{minted}
\end{codeblock}

\section{Conda --- Environnements avec d\'ependances syst\`eme}
\label{sec:conda}
\index{Conda}

\textbf{Conda}\index{Conda} g\`ere non seulement les packages Python mais
aussi les biblioth\`eques syst\`eme (CUDA, MKL, OpenSSL). C'est l'outil
privil\'egi\'e pour les projets ML n\'ecessitant des GPU.

\begin{codeblock}[title=\textbf{Commandes Conda essentielles}]
\begin{minted}{bash}
# Creer un environnement
conda create -n ml-project python=3.11

# Activer
conda activate ml-project

# Installer des packages
conda install numpy pandas scikit-learn
conda install -c pytorch pytorch torchvision  # canal PyTorch

# Exporter l'environnement
conda env export > environment.yml

# Exporter sans les dependances specifiques a l'OS
conda env export --from-history > environment.yml

# Recreer l'environnement
conda env create -f environment.yml

# Lister les environnements
conda env list

# Supprimer un environnement
conda env remove -n ml-project
\end{minted}
\end{codeblock}

\begin{codeblock}[title=\textbf{environment.yml}]
\begin{minted}{yaml}
name: ml-project
channels:
  - pytorch
  - conda-forge
  - defaults
dependencies:
  - python=3.11
  - numpy=1.26
  - pandas=2.1
  - scikit-learn=1.3
  - pytorch=2.1
  - pip:
    - mlflow>=2.8      # packages pip si non dispo en conda
    - wandb>=0.16
\end{minted}
\end{codeblock}

\begin{bonnepratique}[title=\textbf{Conda vs pip}]
\begin{itemize}
  \item Utilisez Conda pour les d\'ependances syst\`eme (CUDA, compilateurs).
  \item Utilisez pip \`a l'int\'erieur d'un environnement Conda pour les
    packages Python purs.
  \item Ne m\'elangez \textbf{jamais} \cli{conda install} et \cli{pip install}
    pour le m\^eme package --- Conda ne verra pas les packages install\'es
    par pip et vice versa.
\end{itemize}
\end{bonnepratique}

\section{Poetry --- Gestion moderne des d\'ependances}
\label{sec:poetry}
\index{Poetry}

\textbf{Poetry}\index{Poetry} est un gestionnaire de d\'ependances moderne
qui combine la gestion des packages, le verrouillage des versions et la
publication de packages dans un seul outil.

\begin{codeblock}[title=\textbf{Workflow Poetry}]
\begin{minted}{bash}
# Installer Poetry
curl -sSL https://install.python-poetry.org | python3 -

# Creer un nouveau projet
poetry new ml-project
cd ml-project

# Ajouter des dependances
poetry add numpy pandas scikit-learn
poetry add --group dev pytest black mypy

# Installer les dependances
poetry install

# Executer dans l'environnement
poetry run python train.py
poetry run pytest tests/

# Exporter en requirements.txt (compatibilite)
poetry export -f requirements.txt -o requirements.txt
\end{minted}
\end{codeblock}

\begin{codeblock}[title=\textbf{pyproject.toml (extrait)}]
\begin{minted}{toml}
[tool.poetry]
name = "ml-project"
version = "0.1.0"
description = "ML classification project"

[tool.poetry.dependencies]
python = "^3.10"
numpy = "^1.26"
pandas = "^2.1"
scikit-learn = "^1.3"

[tool.poetry.group.dev.dependencies]
pytest = "^7.4"
black = "^23.10"
mypy = "^1.6"
\end{minted}
\end{codeblock}

\section{Comparaison des outils}
\label{sec:env-comparison}

\begin{center}
\begin{tabular}{@{}lcccc@{}}
\toprule
\textbf{Crit\`ere} & \textbf{venv+pip} & \textbf{Conda} & \textbf{Poetry} & \textbf{uv} \\
\midrule
Int\'egr\'e Python & Oui & Non & Non & Non \\
D\'ep.\ syst\`eme & Non & Oui & Non & Non \\
Verrouillage & pip-tools & Oui & Oui & Oui \\
GPU/CUDA & Manuel & Oui & Manuel & Manuel \\
Rapidit\'e & Rapide & Lent & Moyen & Tr\`es rapide \\
Publication pkg & Non & Non & Oui & Non \\
Courbe d'apprentissage & Faible & Moyenne & Moyenne & Faible \\
\bottomrule
\end{tabular}
\end{center}

\section{uv --- Le nouvel outil rapide}
\label{sec:uv}
\index{uv}

\textbf{uv} est un gestionnaire de packages Python \'ecrit en Rust,
compatible avec pip, extr\^emement rapide (10--100$\times$ plus rapide
que pip).

\begin{codeblock}[title=\textbf{Workflow uv}]
\begin{minted}{bash}
# Installer uv
curl -LsSf https://astral.sh/uv/install.sh | sh

# Creer un environnement et installer
uv venv
source .venv/bin/activate
uv pip install numpy pandas scikit-learn

# Ou tout en un
uv pip install -r requirements.txt

# Verrouillage
uv pip compile requirements.in -o requirements.txt
uv pip sync requirements.txt
\end{minted}
\end{codeblock}

\section{Bonnes pratiques pour la reproductibilit\'e}
\label{sec:env-best-practices}

\begin{bonnepratique}[title=\textbf{R\`egles d'or pour les environnements}]
\begin{enumerate}
  \item \textbf{Un projet = un environnement}~: ne jamais partager
    d'environnement entre projets.
  \item \textbf{Fichier de d\'ependances versionn\'e}~: \file{requirements.txt}
    ou \file{environment.yml} dans le d\'ep\^ot Git.
  \item \textbf{\'Epingler les versions}~: \cli{numpy==1.26.2}, pas
    \cli{numpy} sans version.
  \item \textbf{S\'eparer dev et prod}~: les outils de d\'eveloppement
    (pytest, black) ne doivent pas \^etre en production.
  \item \textbf{Documenter la version de Python}~: dans le README ou
    \file{pyproject.toml}.
  \item \textbf{Ne jamais modifier l'environnement syst\`eme}~: ne jamais
    utiliser \cli{sudo pip install}.
\end{enumerate}
\end{bonnepratique}

\begin{antipattern}[title=\textbf{Anti-patterns courants}]
\begin{itemize}
  \item Installer tous les packages dans l'environnement syst\`eme
  \item Utiliser \cli{pip install} sans \cli{-r requirements.txt}
  \item Oublier de mettre \`a jour \file{requirements.txt} apr\`es ajout d'un
    package
  \item M\'elanger Conda et pip de mani\`ere anarchique
  \item Partager un environnement Conda via un chemin absolu
\end{itemize}
\end{antipattern}

\section{Mini-projet~: environnement reproductible}
\label{sec:mini-project-ch2}

\begin{miniprojet}[title=\textbf{Mini-projet 2~: Mise en place d'un environnement}]
En reprenant le projet structur\'e du chapitre~\ref{ch:introduction}~:
\begin{enumerate}
  \item Cr\'eez un environnement virtuel avec \cli{venv} ou Conda.
  \item \'Ecrivez un \file{requirements.in} avec les d\'ependances directes.
  \item G\'en\'erez un \file{requirements.txt} verrouill\'e avec \pkg{pip-tools}.
  \item \'Ecrivez un \file{environment.yml} \'equivalent pour Conda.
  \item V\'erifiez que votre coll\`egue peut recr\'eer l'environnement
    \`a partir du fichier de d\'ependances uniquement.
  \item Ajoutez un script \cli{setup.sh} qui automatise la cr\'eation
    de l'environnement.
\end{enumerate}
\end{miniprojet}

\section{Exercices}
\label{sec:exercises-ch2}

\begin{exercice}[$\bigstar$ --- Cr\'eation d'environnement]
Cr\'eez un environnement virtuel avec chacun des trois outils (venv, Conda,
Poetry) pour un projet utilisant \pkg{numpy}, \pkg{pandas} et \pkg{matplotlib}.
Comparez les tailles des environnements et les temps de cr\'eation.
\end{exercice}

\begin{exercice}[$\bigstar\bigstar$ --- R\'esolution de conflits]
Un projet n\'ecessite \pkg{tensorflow>=2.15} et \pkg{numpy<1.24} (conflit
r\'eel). Utilisez \pkg{pip-tools} pour identifier le conflit, puis trouvez
des versions compatibles. Documentez le processus de r\'esolution.
\end{exercice}

\begin{exercice}[$\bigstar\bigstar\bigstar$ --- Environnement multi-plateforme]
Cr\'eez un projet ML qui fonctionne sur Linux, macOS et Windows~:
\begin{enumerate}
  \item \'Ecrivez un \file{environment.yml} Conda compatible multi-OS
    (\cli{conda env export --from-history})
  \item Cr\'eez un script \cli{setup.sh} / \cli{setup.bat} qui d\'etecte
    l'OS et installe les d\'ependances adapt\'ees
  \item Testez sur au moins deux syst\`emes (ou utilisez des conteneurs
    Docker pour simuler)
  \item G\'erez le cas GPU (installer PyTorch CPU vs CUDA selon la machine)
\end{enumerate}
\end{exercice}

\section{Aide-m\'emoire}

\begin{cheatsheet}[title=\textbf{R\'esum\'e du chapitre 2}]
\begin{tabular}{@{}p{4cm}p{8.5cm}@{}}
\toprule
\textbf{Outil} & \textbf{Commandes cl\'es} \\
\midrule
\textbf{venv} & \cli{python -m venv .venv}, \cli{source .venv/bin/activate} \\
\textbf{pip-tools} & \cli{pip-compile requirements.in}, \cli{pip-sync} \\
\textbf{Conda} & \cli{conda create -n env python=3.11}, \cli{conda env export} \\
\textbf{Poetry} & \cli{poetry add pkg}, \cli{poetry install}, \cli{poetry run} \\
\textbf{uv} & \cli{uv pip install}, \cli{uv pip compile}, \cli{uv pip sync} \\
\bottomrule
\end{tabular}
\end{cheatsheet}
